\documentclass{article}
\usepackage[utf8]{inputenc}
\usepackage{graphicx}
\usepackage{grffile}
\usepackage{amsmath}
\usepackage{amssymb}
\usepackage{xcolor}
\usepackage{bm}

\title{Modeling of Supersonic Radiative Marshak Waves Using Simple Models and Advanced Simulations}
\author{Yuval Kochavi}
\date{December 2025}

\begin{document}

\maketitle

\section{Effective mean free path}
if the Rossland mean free path is larger than the halraum diameter, we need to correct it to avoid unphysical results. In other words,
\begin{itemize}
    \item Optically think (diffusion regime): $\lambda_{\mathrm{R}} \ll R$, no correction needed.
    \item Optically thin (transport regime): $\lambda_{\mathrm{R}} \gg R$, need to correct $\lambda_{\mathrm{R}}$ to be $\sim R$.
\end{itemize}
So, we can define an effective mean free path as below:
\begin{equation}
\lambda_{\mathrm{eff}} = \left( \frac{1}{\lambda_{\mathrm{R}}} + \frac{1}{R} \right)^{-1}.
\end{equation}
But, this definition has a problem: when $\lambda_{\mathrm{R}} \ll R$, $\lambda_{\mathrm{eff}} \approx \lambda_{\mathrm{R}}$ (which is good), and when $\lambda_{\mathrm{R}} \gg R$, $\lambda_{\mathrm{eff}} \approx R$ (which is also good). However, in the intermediate regime, $\lambda_{\mathrm{eff}}$ can be significantly smaller than both $\lambda_{\mathrm{R}}$ and $R$, which is unphysical. To fix this, we can use the following definition:
\begin{equation}
\lambda_{\mathrm{eff}} = \left( \frac{1}{\lambda_{\mathrm{R}}^n} + \frac{1}{R^n} \right)^{-\frac{1}{n}} \to \min\left( \lambda_{\mathrm{R}}, R \right).
\end{equation}
This way, $\lambda_{\mathrm{eff}}$ smoothly transitions between $\lambda_{\mathrm{R}}$ and $R$ without becoming unphysically small in the intermediate regime. The choice of $n$ controls the sharpness of the transition; a larger $n$ results in a sharper transition.

\subsection*{Effective mean-free-path correction in the numerical model}

To prevent unphysical behavior when the Rosseland mean free path becomes comparable to or larger than the hohlraum radius, an effective mean-free-path correction is implemented in the numerical model.  
The Rosseland mean free path is evaluated as
\begin{equation}
\lambda_{\mathrm{R}}(T_s,\rho)
=
g\,T_s^{\alpha}\,\rho^{-(\lambda+1)},
\end{equation}
and the geometric transport scale is taken as the hohlraum diameter,
\begin{equation}
\lambda_{\mathrm{geom}} = 2R.
\end{equation}

A smooth minimum between these two scales is then defined using a generalized power mean,
\begin{equation}
\lambda_{\mathrm{eff}}
=
\left(
\lambda_{\mathrm{geom}}^{-p}
+
\lambda_{\mathrm{R}}^{-p}
\right)^{-1/p},
\end{equation}
where the numerical parameter $p$ controls the sharpness of the transition.  
This construction ensures that
\[
\lambda_{\mathrm{eff}} \to
\begin{cases}
\lambda_{\mathrm{R}}, & \lambda_{\mathrm{R}} \ll 2R,\\
2R, & \lambda_{\mathrm{R}} \gg 2R,
\end{cases}
\]
without introducing an artificial minimum at intermediate optical depth.

The opacity prefactor is renormalized consistently with the effective mean free path,
\begin{equation}
g_{\mathrm{eff}}(t)
=
g\,\frac{\lambda_{\mathrm{eff}}(t)}{\lambda_{\mathrm{R}}(t)},
\end{equation}
so that the diffusion coefficient scales as $D \propto \lambda_{\mathrm{eff}}$ in both the diffusion and geometry-limited regimes.

When enabled (\texttt{lam\_eff=True}), this correction is applied dynamically at each time step using the instantaneous surface temperature $T_s(t)$ (or $T_s(t-\Delta t)$ in the Marshak marcher).  
In the ablation case with density evolution, the correction is coupled directly to the density-dependent transport scaling by integrating the combined factor
\[
\left(\tfrac{g_{\mathrm{eff}}}{g}\right)\rho^{\,\mu-\lambda-2}
\]
inside the time integral used to construct the similarity coefficient, ensuring that geometry-limited transport and density evolution are treated consistently rather than as separable corrections.

Finally, the dimensionless similarity parameter controlling the Marshak-wave front evolution is evaluated as
\begin{equation}
Z_1 = Z_1\!\left(C_{\mathrm{eff}}, \rho\right),
\end{equation}
so that the wave-front dynamics, opacity, geometry, and density are consistently coupled through $\lambda_{\mathrm{eff}}$.

\bigskip
\noindent
\textbf{Physical meaning:}  
This construction enforces a physically consistent transition between diffusion-dominated transport and free-streaming transport, while preserving correct scaling of opacity, diffusion coefficient, and Marshak-wave similarity structure in cylindrical hohlraum geometry.

\section{Results}
\begin{table}[h!]
\centering
\begin{tabular}{ccccccccc}
\hline
$Diameter$ & L & $\frac{L}{\lambda_R}$ & $\frac{\lambda_R}{D}$ & $\rho$ & $\lambda_R\;[\mathrm{cm}]$ & $n$ & Material & Expt. \\
\hline
0.06 & 0.04 & 0.11 & 5.93 & 0.05   & 0.356 & $\to\infty$ & C$_6$H$_{12}$ & Xu \\
0.06 & 0.04 & 0.41 & 1.633 & 0.05   & 0.098 & $>2$        & C$_6$H$_{12}$Cu$_{0.394}$ & Xu \\
0.16 & 0.15 & 8.33 & 0.1125 & 0.04   & 0.018 & $\ge 2$     & Ta$_2$O$_5$ & Back \\
0.16 & 0.15 & 0.83 & 1.125 & 0.05   & 0.18 & 1.5     & SiO$_2$ & Back \\
0.30 & 0.15 & 1.36 & 0.3667 & 0.01   & 0.11 & $\ge 2$     & SiO$_2$ low energy & Back \\
0.20 & 0.28 & 1.32 & 1.06 & 0.1249 & 0.212 & 1           & SiO$_2$ & Moore \\
0.20 & 0.28 & 0.28 & 4.94 & 0.1139 & 0.988 & 2--2.5      & C$_8$H$_7$Cl & Moore \\
0.08 & 0.12 & 0.31 & 4.875 & 0.0065  & 0.39  & 1           & C$_{15}$H$_{20}$O$_6$ & Keiter \\
0.08 & 0.12 & 2.11 & 0.7125 & 0.0625 & 0.057 & 2           & C$_{15}$H$_{20}$O$_6$Au$_{0.172}$ & Keiter \\
0.02 & 0.03 & 0.64 & 2.35 & 0.16   & 0.047 & ? ($n>3$)   & C$_8$H$_8$ & Ji-Yan \\
0.10 & 0.12 & 0.29 & 4.1 & 0.029  & 0.41  & 1           & French gold & French3 \\
0.20 & 0.2 & 0.23 & 4.33 & 0.0189  & 0.86  & 1           & French copper & French3 \\
\hline
\end{tabular}
\caption{Parameters extracted from maximum temperature mean free path calculation.}
\label{tab:macimu depth}
\end{table}


\begin{table}[h!]
\centering
\begin{tabular}{c c c c c c c c c}
$Diameter$ & L & $\frac{L}{\lambda_R}$ & $\frac{\lambda_R}{D}$ & $\rho$ & $\lambda_R\;[\mathrm{cm}]$ & $n$ & Material & Expt. \\
\hline
0.06 & 0.04 & 0.381 & 1.75 & 0.05   & 0.105 & $\to\infty$ & C$_6$H$_{12}$ & Xu \\
0.06 & 0.04 & 1.33 & 0.50 & 0.05   & 0.03 & $>2$        & C$_6$H$_{12}$Cu$_{0.394}$ & Xu \\
0.16 & 0.15 & 10.0 & 0.0938 & 0.04   & 0.015 & $\ge 2$     & Ta$_2$O$_5$ & Back \\
0.16 & 0.15 & 1.88 & 0.50 & 0.05   & 0.08 & 1.5     & SiO$_2$ & Back \\
0.30 & 0.15 & 3.00 & 0.1667 & 0.01   & 0.05 & $\to\infty$     & SiO$_2$ low energy & Back \\
0.20 & 0.28 & 1.32 & 1.06 & 0.1249 & {\color{red}0.212} & 1           & SiO$_2$ & Moore \\
0.20 & 0.28 & 1.22 & 1.15 & 0.1139 & 0.23 & 2--2.5      & C$_8$H$_7$Cl & Moore \\
0.08 & 0.12 & 1.46 & 1.025 & 0.0065  & 0.082  & 1           & C$_{15}$H$_{20}$O$_6$ & Keiter \\
0.08 & 0.12 & 4.00 & 0.375 & 0.0625 & 0.03 & 2           & C$_{15}$H$_{20}$O$_6$Au$_{0.172}$ & Keiter \\
0.02 & 0.03 & 2.00 & 0.75 & 0.16   & 0.015 & ? ($n>3$)   & C$_8$H$_8$ & Ji-Yan \\
0.10 & 0.12 & 1.14 & 1.05 & 0.029  & 0.105  & 1           & French gold & French3 \\
0.20 & 0.2 & 1.33 & 0.75 & 0.0189  & 0.15  & 1           & French copper & French3 \\
\hline
\end{tabular}
\caption{Parameters extracted from optical depth calculation}
\label{tab:optical depth}
\end{table}
\end{document}